\documentclass[11pt]{amsart}
\usepackage{calc,amssymb,amsmath,ulem,stmaryrd,fullpage}
\usepackage{enumerate}
\usepackage{alltt}
\RequirePackage[dvipsnames,usenames]{color}
\let\mffont=\sf
\normalem
%\input{kmacros3.sty}
\input{xy}
\xyoption{all}
\usepackage{hyperref}
\usepackage{graphicx}
\usepackage[all,cmtip]{xy}
\usepackage{verbatim}
\usepackage[left=0.95in,top=0.95in,right=0.95in,bottom=0.95in]{geometry}
\newcommand{\tauCohomology}{T}
\newcommand{\FCohomology}{S}


\theoremstyle{theorem}
%\newtheorem*{mainthm*}{Main Theorem}

\newcommand{\note}[1]{\marginpar{\sffamily\tiny #1}}

\def\todo#1{\textcolor{Mahogany}%
{\sffamily\footnotesize\newline{\color{Mahogany}\fbox{\parbox{\textwidth-15pt}{#1}}}\newline}}

\renewcommand{\O}{\mathcal O}
\newcommand{\ie}{{\itshape i.e.} }
\newcommand{\roundup}[1]{\lceil #1 \rceil}
\newcommand{\rounddown}[1]{\lfloor #1 \rfloor}
\renewcommand{\to}[1][]{\xrightarrow{\ #1\ }}
\newcommand{\onto}[1][]{\protect{\xrightarrow{\ #1\ }\hspace{-0.8em}\rightarrow}}
\newcommand{\into}[1][]{\lhook \joinrel \xrightarrow{\ #1\ }}
%\newcommand{DBDual}[1]{{\underline{\omega}_{#1}}}
\newcommand{\DBDual}[1]{{{\uuline \omega}{}^{\mydot}_{#1}}}
\newcommand{\Tt}{{\mathfrak{T}}}
%\newcommand{\bn}{{\mathfrak{n}} }
\newcommand{\sol}[1]{ \vskip 6pt{\color{blue} {\bf Solution:} #1} \vskip 6pt}

\newcommand{\bR}{{\mathbb{R}}}
\newcommand{\va}{{\vec{a}}}
\newcommand{\vb}{{\vec{b}}}
\newcommand{\vzero}{{\vec{0}}}
\newcommand{\vi}{{\vec{i}}}
\newcommand{\vj}{{\vec{j}}}
\newcommand{\vk}{{\vec{k}}}
\newcommand{\vh}{{\vec{h}}}
\newcommand{\vr}{{\vec{r}}}
\newcommand{\vu}{{\vec{u}}}
\newcommand{\vv}{{\vec{v}}}
\newcommand{\vw}{{\vec{w}}}
\newcommand{\vx}{{\vec{x}}}
\newcommand{\vy}{{\vec{y}}}
\newcommand{\vz}{{\vec{z}}}
\newcommand{\vT}{{\vec{T}}}
\newcommand{\vN}{{\vec{N}}}
\newcommand{\vF}{{\vec{F}}}
\newcommand{\vG}{{\vec{G}}}
\newcommand{\bF}{\mathbb{F}}
\newcommand{\bQ}{\mathbb{Q}}
\newcommand{\bZ}{\mathbb{Z}}
\newcommand{\bC}{\mathbb{C}}
\newcommand{\Gal}{\textnormal{Gal}}
\newcommand{\Aut}{\textnormal{Aut}}
\newcommand{\Emb}{\textnormal{Emb}}
\newcommand{\Hom}{\textnormal{Hom}}
\newcommand{\Ext}{\textnormal{Ext}}
\newcommand{\Tor}{\textnormal{Tor}}

\newcommand{\coker}{\textnormal{coker}}
\newcommand{\id}{\textnormal{id}}


\begin{document}
\title{WS \#7 -- Math 6310\\Fall 2019}
\author{Due: Wednesday December 4th}
\maketitle

You are encouraged to work in groups on this.  Only one write up needs to be turned in per group.  
\vskip 6pt
\noindent
{\bf 1.}  Suppose that $P_{\bullet}, P'_{\bullet}$ are projective resolutions of $R$-modules $M$ and $N$.  Suppose have an $R$-module map $h : M \to N$.  Consider the diagram
\[
\xymatrix{
\ldots \ar[d]  & \ldots \ar[d] \\
P_{2} \ar[d]_d \ar@{.>}[r] & P'_{2} \ar[d]^{d'} \\
P_{1} \ar[d] \ar@{.>}[r] & P'_{1} \ar[d] \\
P_{0} \ar@{->>}[d] \ar@{.>}[r] & P'_{0} \ar@{->>}[d] \\
M \ar[r]_{h} & N \\
}
\]
Show that there exist the dotted arrows making the diagram commute.  This gives us a map of chain complexes $\phi : P_{\bullet} \to P'_{\bullet}$.  Was the map you constructed unique?
\newpage 
\noindent
{\bf Definition:}  Two maps of chain complexes $\phi, \psi: K_{\bullet} \to L_{\bullet}$ are declared to be \emph{homotopic} (denoted $\phi \sim \psi$ frequently) if for each $n$, there is an $R$-module map $h_n : K_n \to L_{n+1}$ such that
\[
\phi_n - \psi_n = d_L h_n + h_{n-1} d_K
\]
\vskip 9pt
\noindent
{\bf 2.}  If $\phi, \psi$ are homotopic, show that they induce the same maps on homology.  
\vskip 3pt
\emph{Hint:}  First draw a diagram picturing the above equation.
\newpage
\noindent
{\bf 3.}  Back in the setting of {\bf 1.}, show that any two maps $\phi, \psi : P_{\bullet} \to P'_{\bullet}$ you constructed are homotopic.  
\newpage
\noindent
{\bf 4.}  
Conclude that this homotopy is preserved after tensoring by another module $A$ and thus conclude that $\mathrm{Tor}_i(M, A)$ is independent of the choice of projective resolution of $M$.  Likewise conclude that $\mathrm{Ext}^i(M, A)$ is independent of the projective resolution of $M$.
\newpage
\noindent
{\bf 5.}  Compute a free resolution of the ideal $\langle x,y \rangle$ inside $R = k[x,y]$.
\newpage
\noindent
{\bf 6.}  Use this to compute directly by hand $\Tor_i( \langle x,y \rangle, k[x,y]/\langle x,y \rangle)$.
\newpage
\noindent
{\bf 7.}  Likewise compute directly $\Ext^i( \langle x,y \rangle, k[x,y]/\langle x,y \rangle)$.

\end{document}

